\makeatletter
\def\rap@fontpath{../../../Common/}
\makeatother
\documentclass{../../../Common/Latex/Templates/rapport}
% ---------------------------------------------------------------------------
% Shared settings for the H2026 General Assembly document set.
%
% Usage, directly after \documentclass{../../../Common/Latex/Templates/rapport}
% in a document one folder below this file:
%   % ---------------------------------------------------------------------------
% Shared settings for the H2026 General Assembly document set.
%
% Usage, directly after \documentclass{../../../Common/Latex/Templates/rapport}
% in a document one folder below this file:
%   % ---------------------------------------------------------------------------
% Shared settings for the H2026 General Assembly document set.
%
% Usage, directly after \documentclass{../../../Common/Latex/Templates/rapport}
% in a document one folder below this file:
%   \input{../ga-common}
%
% Image paths are relative to the compiled document's folder, so every
% document using this file must sit one folder below it.
% ---------------------------------------------------------------------------

% Meeting details, used on covers and in each document's ID block.
\newcommand{\meetingdate}{30 September 2026}
\newcommand{\meetingtime}{12:00}
\newcommand{\meetinglocation}{Festsalen, NMBU}
\newcommand{\meetingline}{\textbf{Meeting:} \meetingdate, \meetingtime, \meetinglocation}

% Header: moon logo on the left, full-text logo on the right. \projectname is
% still set because the cover title uses it.
\projectname{Eclipse NMBU}
\projectlogo{../../../Common/Eclipse Images/EclipseMoon}
\orglogo{../../../Common/Eclipse Images/EclipseFullText}

%
% Image paths are relative to the compiled document's folder, so every
% document using this file must sit one folder below it.
% ---------------------------------------------------------------------------

% Meeting details, used on covers and in each document's ID block.
\newcommand{\meetingdate}{30 September 2026}
\newcommand{\meetingtime}{12:00}
\newcommand{\meetinglocation}{Festsalen, NMBU}
\newcommand{\meetingline}{\textbf{Meeting:} \meetingdate, \meetingtime, \meetinglocation}

% Header: moon logo on the left, full-text logo on the right. \projectname is
% still set because the cover title uses it.
\projectname{Eclipse NMBU}
\projectlogo{../../../Common/Eclipse Images/EclipseMoon}
\orglogo{../../../Common/Eclipse Images/EclipseFullText}

%
% Image paths are relative to the compiled document's folder, so every
% document using this file must sit one folder below it.
% ---------------------------------------------------------------------------

% Meeting details, used on covers and in each document's ID block.
\newcommand{\meetingdate}{30 September 2026}
\newcommand{\meetingtime}{12:00}
\newcommand{\meetinglocation}{Festsalen, NMBU}
\newcommand{\meetingline}{\textbf{Meeting:} \meetingdate, \meetingtime, \meetinglocation}

% Header: moon logo on the left, full-text logo on the right. \projectname is
% still set because the cover title uses it.
\projectname{Eclipse NMBU}
\projectlogo{../../../Common/Eclipse Images/EclipseMoon}
\orglogo{../../../Common/Eclipse Images/EclipseFullText}

\usepackage{enumitem}

\reporttitle{GF for Dummies}

% ---------------------------------------------------------------------------
% Cover / title page for rapport-class documents.
%
% Usage (after \documentclass{rapport} and before \begin{document}):
%   % ---------------------------------------------------------------------------
% Cover / title page for rapport-class documents.
%
% Usage (after \documentclass{rapport} and before \begin{document}):
%   % ---------------------------------------------------------------------------
% Cover / title page for rapport-class documents.
%
% Usage (after \documentclass{rapport} and before \begin{document}):
%   \input{../cover/rapport-cover}
%   \missionpatch{path/to/patch}      % optional circular mission/team patch;
%                                     % without it, the org logo is shown large
%                                     % at the top instead
%   \coverkind{Vedtekter}             % optional small line above the name
%   \reportsubtitle{Optional subtitle line}
%   \orgname{Eclipse NMBU}{Norwegian University of Life Sciences, Norway}
%   \reportlocation{Ås, Norge}
%   \date{3. oktober 2022}            % optional, defaults to \today
%   \contactinfo{Contact: post@eclipsenmbu.no}  % optional, shown at the bottom
%
% Then, as the first thing inside \begin{document}:
%   \makecoverpage
%
% The large title line shows \projectname (falling back to \reporttitle if
% \coverkind is not set), so a document typically sets both:
%   \projectname{Eclipse NMBU}
%   \coverkind{Vedtekter}
% ---------------------------------------------------------------------------

\makeatletter
\def\@missionpatch{}
\newcommand{\missionpatch}[1]{\def\@missionpatch{#1}}

\def\@coverkind{}
\newcommand{\coverkind}[1]{\def\@coverkind{#1}}

\def\@reportsubtitle{}
\newcommand{\reportsubtitle}[1]{\def\@reportsubtitle{#1}}

\def\@orgnameone{}
\def\@orgnametwo{}
\newcommand{\orgname}[2]{\def\@orgnameone{#1}\def\@orgnametwo{#2}}

\def\@reportlocation{}
\newcommand{\reportlocation}[1]{\def\@reportlocation{#1}}

\def\@contactinfo{}
\newcommand{\contactinfo}[1]{\def\@contactinfo{#1}}

\newcommand{\makecoverpage}{%
  \begingroup
  \thispagestyle{empty}
  \begin{center}
  \ifx\@missionpatch\@empty
    \includegraphics[width=10cm]{\@orglogo}\par
  \else
    \includegraphics[width=6.5cm]{\@missionpatch}\par
  \fi
  \vspace{3.4cm}
  \ifx\@coverkind\@empty
    {\sffamily\Huge\bfseries \@reporttitle \par}
  \else
    {\sffamily\Huge\bfseries \@coverkind \par}
    {\sffamily\Huge\bfseries \@projectname \par}
  \fi
  \ifx\@reportsubtitle\@empty\else
    \vspace{1.6cm}
    {\sffamily\bfseries \@reportsubtitle \par}
  \fi
  \vspace{2cm}
  \ifx\@missionpatch\@empty\else
    \includegraphics[height=1.3cm]{\@orglogo}\par
    \vspace{0.6cm}
  \fi
  \@orgnameone \\
  \@orgnametwo \par
  \vspace{0.6cm}
  \ifx\@reportlocation\@empty
    \@date
  \else
    \@reportlocation, \@date
  \fi
  \par
  \vfill
  \ifx\@contactinfo\@empty\else
    {\small \@contactinfo \par}
  \fi
  \end{center}
  \endgroup
  \newpage
  \pagenumbering{roman}
  \pagestyle{plain}
}
\makeatother

%   \missionpatch{path/to/patch}      % optional circular mission/team patch;
%                                     % without it, the org logo is shown large
%                                     % at the top instead
%   \coverkind{Vedtekter}             % optional small line above the name
%   \reportsubtitle{Optional subtitle line}
%   \orgname{Eclipse NMBU}{Norwegian University of Life Sciences, Norway}
%   \reportlocation{Ås, Norge}
%   \date{3. oktober 2022}            % optional, defaults to \today
%   \contactinfo{Contact: post@eclipsenmbu.no}  % optional, shown at the bottom
%
% Then, as the first thing inside \begin{document}:
%   \makecoverpage
%
% The large title line shows \projectname (falling back to \reporttitle if
% \coverkind is not set), so a document typically sets both:
%   \projectname{Eclipse NMBU}
%   \coverkind{Vedtekter}
% ---------------------------------------------------------------------------

\makeatletter
\def\@missionpatch{}
\newcommand{\missionpatch}[1]{\def\@missionpatch{#1}}

\def\@coverkind{}
\newcommand{\coverkind}[1]{\def\@coverkind{#1}}

\def\@reportsubtitle{}
\newcommand{\reportsubtitle}[1]{\def\@reportsubtitle{#1}}

\def\@orgnameone{}
\def\@orgnametwo{}
\newcommand{\orgname}[2]{\def\@orgnameone{#1}\def\@orgnametwo{#2}}

\def\@reportlocation{}
\newcommand{\reportlocation}[1]{\def\@reportlocation{#1}}

\def\@contactinfo{}
\newcommand{\contactinfo}[1]{\def\@contactinfo{#1}}

\newcommand{\makecoverpage}{%
  \begingroup
  \thispagestyle{empty}
  \begin{center}
  \ifx\@missionpatch\@empty
    \includegraphics[width=10cm]{\@orglogo}\par
  \else
    \includegraphics[width=6.5cm]{\@missionpatch}\par
  \fi
  \vspace{3.4cm}
  \ifx\@coverkind\@empty
    {\sffamily\Huge\bfseries \@reporttitle \par}
  \else
    {\sffamily\Huge\bfseries \@coverkind \par}
    {\sffamily\Huge\bfseries \@projectname \par}
  \fi
  \ifx\@reportsubtitle\@empty\else
    \vspace{1.6cm}
    {\sffamily\bfseries \@reportsubtitle \par}
  \fi
  \vspace{2cm}
  \ifx\@missionpatch\@empty\else
    \includegraphics[height=1.3cm]{\@orglogo}\par
    \vspace{0.6cm}
  \fi
  \@orgnameone \\
  \@orgnametwo \par
  \vspace{0.6cm}
  \ifx\@reportlocation\@empty
    \@date
  \else
    \@reportlocation, \@date
  \fi
  \par
  \vfill
  \ifx\@contactinfo\@empty\else
    {\small \@contactinfo \par}
  \fi
  \end{center}
  \endgroup
  \newpage
  \pagenumbering{roman}
  \pagestyle{plain}
}
\makeatother

%   \missionpatch{path/to/patch}      % optional circular mission/team patch;
%                                     % without it, the org logo is shown large
%                                     % at the top instead
%   \coverkind{Vedtekter}             % optional small line above the name
%   \reportsubtitle{Optional subtitle line}
%   \orgname{Eclipse NMBU}{Norwegian University of Life Sciences, Norway}
%   \reportlocation{Ås, Norge}
%   \date{3. oktober 2022}            % optional, defaults to \today
%   \contactinfo{Contact: post@eclipsenmbu.no}  % optional, shown at the bottom
%
% Then, as the first thing inside \begin{document}:
%   \makecoverpage
%
% The large title line shows \projectname (falling back to \reporttitle if
% \coverkind is not set), so a document typically sets both:
%   \projectname{Eclipse NMBU}
%   \coverkind{Vedtekter}
% ---------------------------------------------------------------------------

\makeatletter
\def\@missionpatch{}
\newcommand{\missionpatch}[1]{\def\@missionpatch{#1}}

\def\@coverkind{}
\newcommand{\coverkind}[1]{\def\@coverkind{#1}}

\def\@reportsubtitle{}
\newcommand{\reportsubtitle}[1]{\def\@reportsubtitle{#1}}

\def\@orgnameone{}
\def\@orgnametwo{}
\newcommand{\orgname}[2]{\def\@orgnameone{#1}\def\@orgnametwo{#2}}

\def\@reportlocation{}
\newcommand{\reportlocation}[1]{\def\@reportlocation{#1}}

\def\@contactinfo{}
\newcommand{\contactinfo}[1]{\def\@contactinfo{#1}}

\newcommand{\makecoverpage}{%
  \begingroup
  \thispagestyle{empty}
  \begin{center}
  \ifx\@missionpatch\@empty
    \includegraphics[width=10cm]{\@orglogo}\par
  \else
    \includegraphics[width=6.5cm]{\@missionpatch}\par
  \fi
  \vspace{3.4cm}
  \ifx\@coverkind\@empty
    {\sffamily\Huge\bfseries \@reporttitle \par}
  \else
    {\sffamily\Huge\bfseries \@coverkind \par}
    {\sffamily\Huge\bfseries \@projectname \par}
  \fi
  \ifx\@reportsubtitle\@empty\else
    \vspace{1.6cm}
    {\sffamily\bfseries \@reportsubtitle \par}
  \fi
  \vspace{2cm}
  \ifx\@missionpatch\@empty\else
    \includegraphics[height=1.3cm]{\@orglogo}\par
    \vspace{0.6cm}
  \fi
  \@orgnameone \\
  \@orgnametwo \par
  \vspace{0.6cm}
  \ifx\@reportlocation\@empty
    \@date
  \else
    \@reportlocation, \@date
  \fi
  \par
  \vfill
  \ifx\@contactinfo\@empty\else
    {\small \@contactinfo \par}
  \fi
  \end{center}
  \endgroup
  \newpage
  \pagenumbering{roman}
  \pagestyle{plain}
}
\makeatother

\coverkind{General Assembly for Dummies (GF for Dummies)}
\reportsubtitle{A plain-language, non-binding guide -- 26H-ECL-GA-DUMMY1}
\orgname{Eclipse NMBU}{Norwegian University of Life Sciences, Norway}
\reportlocation{\meetinglocation, Ås}
\date{\meetingdate, \meetingtime}
\contactinfo{Contact: post@eclipsenmbu.no}

\begin{document}

\makecoverpage
\tableofcontents
\reportmainmatter

\noindent\textbf{Document ID:} 26H-ECL-GA-DUMMY1\\
\meetingline\\
\textbf{This document is not binding.} If anything here conflicts with the Bylaws, the Notice of Meeting, the Agenda, or any other formal document in this set, those documents control -- this is just a friendly explanation to help you follow along at your first General Assembly (generalforsamling, or ``GF'').

\vsection{What is a General Assembly?}
The General Assembly (generalforsamling) is Eclipse NMBU's supreme decision-making meeting. It's where the whole membership -- not just the Board -- gets to vote on the big things: the budget, the rules (Bylaws), and who sits on the Board next.

\vsection{Who can vote?}
Board Members, Core Members and Active Members can vote. Everyone -- including Inactive Members -- is welcome to attend, listen, and speak, but only those three tiers can actually cast a vote or stand for election. See the Glossary (26H-ECL-GA-GLOSS1) if you're not sure which tier you're in.

\vsection{Constitution of the Meeting}
Before anything else can happen, the meeting has to ``constitute'' itself: elect someone to chair the meeting (Meeting Chair), someone to take minutes (Secretary), two people to count votes (Vote Counters), and two people to check and sign the final minutes (Protocol Signers), record how many eligible voters are present, and approve the Notice and the Agenda. All of this is a single Case, Case 1 -- there is no separate Case for approving the Notice or the Agenda. Nothing else can be validly decided until this is done.

\vsection{The Agenda}
The Agenda (saksliste) is the running order for the meeting -- the list of everything that will be discussed and voted on, in order. This meeting has eight Cases, numbered 1 to 8. You can find it as 26H-ECL-GA-AGD1, and the Overview (26H-ECL-GA-OVW1) lists every document you might be handed.

\vsection{What's a Case?}
A Case (sak) is one item of business -- one thing the General Assembly is being asked to decide. Each Case has its own Case Paper.

\vsection{How to read a Case Paper}
Every Case Paper (sakspapir) follows the same template (26H-ECL-GA-TPL1), so once you know the shape, you can read any of them: what's the background (why are we talking about this?), what's the actual proposal (what are we being asked to vote for?), and any other information you need to make up your mind.

\vsection{Right to Propose}
Right to Propose (forslagsrett) is the right to submit a new Case, or an Amendment to an existing one. At this General Assembly, that means Board Members, Core Members and Active Members -- the same people who can vote.

\vsection{How to submit an Amendment}
An Amendment (endringsforslag) is a proposed change to a Case that's already on the Agenda. You write down exactly what you want changed and what you want it changed to, and submit it in writing before voting on Cases begins. Full details are in the Amendment Submission Procedure, 26H-ECL-GA-PROC1. \textbf{Once voting on Cases has started, no new Amendments can be submitted} -- there's no last-minute exception.

\vsection{Floor Amendments -- really, Floor Nominations}
You'll sometimes hear ``floor amendment'' (benkeforslag) used loosely, but in this document set it has one specific meaning: nominating a \emph{candidate} for election directly from the floor, during the meeting itself, rather than in advance. It's not a way to change a Case's text -- that's what Amendments are for.

\vsection{Advance Nominations}
If you want to run for a Board position, you can let the Board know before the meeting -- that's an Advance Nomination (forhåndsbenking). You'll then appear as a candidate on the Agenda. You can also be nominated from the floor on the day.

\vsection{Hand signals}
If the Meeting Chair uses hand signals to speed up discussion (e.g. finger-wave agreement, crossed arms for a point of order), they explain them at the start of the meeting.

\vsection{Voting}
Most votes are decided by simple majority -- more than half of the votes cast, not counting blanks. Voting is usually by show of hands, unless someone asks for a written and secret vote.

\vsection{Two-thirds majority}
Some decisions need more than a simple majority. Changing the Bylaws needs \textbf{2/3 of all Eclipse NMBU members} -- not just those in the room -- to vote in favour. That's a high bar on purpose: it's the Association's constitution.

\vsection{Election procedure}
Board elections are one role at a time. If a candidate gets a majority, they're elected. If nobody does, there's a second round between the top candidates. Blank votes count in the first round, but can be squeezed out in a second round if nobody wins.

\vsection{Voting Order}
Before voting starts, the Meeting Chair proposes an order to vote in -- usually Amendments first, then the underlying proposal -- designed to use as few individual votes as possible. See 26H-ECL-GA-VOTORD1 for the full logic, including the special two-stage process for the Bylaw Amendment.

\vsection{Protocol}
The Protocol (protokoll) is the official written record of everything the General Assembly decided -- every proposal, every vote count, every election result. Once the two Protocol Signers have checked and signed it, it's the final, authoritative account of the meeting.

\vsection{When do decisions become effective?}
Most resolutions take effect as soon as they're passed. Bylaw Amendments are a bit different: they take effect once the Protocol is signed, unless the General Assembly says otherwise.

\vsection{What happens after the General Assembly?}
The Protocol gets published (within two weeks). Newly elected Board members take up their roles (with an overlap/handover period). If the Bylaw Amendment passed, work starts on rolling out the new structure -- new Membership Contracts, Responsibility Contracts for anyone taking on a leadership role, and so on.

\vsection{How this all fits together}
It can look like a lot of separate documents. Here is the map:
\begin{itemize}
  \item \textbf{Before the meeting}, the Overview (26H-ECL-GA-OVW1) lists every document in the set and what it is for; the Notice of Meeting (26H-ECL-GA-NOT1) tells you when and where, and what deadlines apply; the Agenda (26H-ECL-GA-AGD1) lists the Cases in order, each pointing to its own Case Paper.
  \item \textbf{To act before or during the meeting}, the Amendment Submission Procedure (26H-ECL-GA-PROC1) is how you change a Case's wording, and Candidate Documentation (26H-ECL-GA-CAND1) is how you stand for election.
  \item \textbf{During the meeting}, the Voting Order (26H-ECL-GA-VOTORD1) and the Full Meeting Procedure (26H-ECL-GA-PROC2) set out how discussion and voting actually run, Case by Case; ballots (where used) and the Meeting Protocol (26H-ECL-GA-PROT1) capture what is actually decided as it happens.
  \item \textbf{The Bylaw Amendment} is its own thread running through several documents at once: the proposed text and Case Paper (26H-ECL-GA-CASE4A), its plain-language explanation (26H-ECL-GA-CASE4B), and its two-stage voting procedure (26H-ECL-GA-CASE4C) -- and its result decides which version of two other Cases (C-suite structure and Election of Board Members) actually gets used.
  \item \textbf{After the meeting}, the signed Protocol is the authoritative record of everything decided. If the Bylaw Amendment passed, the Membership Contract (26H-ECL-GA-MEMB1) and the Responsibility \& Signature Contract (26H-ECL-GA-RESP1B, logged via 26H-ECL-GA-RESP2) are the documents that put the New Bylaws into practice day to day.
  \item \textbf{Throughout}, this guide and the Glossary (26H-ECL-GA-GLOSS1) are the two non-binding documents: everything else in the list above is a binding or procedural document that these two just help you read.
\end{itemize}

\end{document}
